French number words are compositionally systematic but atypical, especially in the 70--99 range (e.g., \textit{soixante-dix}, \textit{quatre-vingt-dix}). We test whether current large language models show measurable weakness on these irregular forms. We introduce \frcountmap, a reproducible black-box probe pipeline that evaluates numeral transduction and pattern-level classification with deterministic decoding. Our synthetic suite covers every integer from 0 to 199, and we complement it with \mgsmfr for broader mathematical reasoning.

Across 1,935 analyzed outputs, both models achieve near-ceiling performance on direct transduction. For \gptfourone, \digitword reaches 98.6\%--100\% across all bands and \worddigit is 100\% in every complete condition. \gptfouronemini remains similarly strong (96.7\%--100\% on \digitword; 100\% on \worddigit). Two-proportion tests show no significant irregular-band penalty for transduction after FDR correction. The strongest band effect appears in \patterncls, where irregular items are easier under our label schema (e.g., risk difference $+0.435$, FDR-corrected $p<0.001$ for \gptfourone direct). On \mgsmfr, reasoning prompts improve \gptfourone from 48.8\% to 57.5\%.

These results show that modern API LLMs can reliably map French numerals under low-temperature prompting, while conclusions about ``difficulty'' depend strongly on probe design.
